% Source: https://artofproblemsolving.com/community/c3h2805128_2010_midmichigan_mathematical_olympiad__grades_79
% Stable identity: aops-topic-2805128-q05
Let the rates be $r_i$ with $S=\sum r_i$. Person $i$ eats for $t_i=\dfrac1{2(S-r_i)}$ and eats $r_it_i$. So
\[1=\sum\dfrac{r_i}{2(S-r_i)},\qquad T=\sum t_i=\sum\dfrac1{2(S-r_i)}.\]
Since $\dfrac{S}{S-r_i}=1+\dfrac{r_i}{S-r_i}$,
\[S\,T=\sum\dfrac{S}{2(S-r_i)}=\sum\dfrac12+\sum\dfrac{r_i}{2(S-r_i)}=2+1=3.\]
Together they need time $\dfrac1S=\dfrac T3$, so they would be $\boxed{3}$ times faster.
